%=============================================================================
% Page_Geilker_Literature.tex
% Literature Review: Semiclassical Gravity, the Page-Geilker Experiment,
% and Constraint-Field Approaches to Quantum Gravity
% Cross-referenced with DFD's Quasi-Static Constraint-Field Picture
% Date: 2026-03-23
%=============================================================================
\documentclass[12pt,a4paper]{article}
\usepackage[utf8]{inputenc}
\usepackage{amsmath,amssymb,amsfonts,amsthm}
\usepackage{hyperref}
\usepackage{booktabs}
\usepackage{longtable}
\usepackage{geometry}
\usepackage{xcolor}
\usepackage{tcolorbox}
\geometry{margin=2.5cm}

\newtheorem{theorem}{Theorem}
\newtheorem{proposition}{Proposition}
\newtheorem{lemma}{Lemma}
\newtheorem{corollary}{Corollary}
\newtheorem{remark}{Remark}

\definecolor{dfdgreen}{RGB}{0,128,0}
\definecolor{dfdblue}{RGB}{0,50,180}
\definecolor{dfdred}{RGB}{200,0,0}
\definecolor{dfdyellow}{RGB}{200,180,0}

\newcommand{\GREEN}{\textcolor{dfdgreen}{\textbf{GREEN}}}
\newcommand{\YELLOW}{\textcolor{dfdyellow}{\textbf{YELLOW}}}
\newcommand{\RED}{\textcolor{dfdred}{\textbf{RED}}}

\title{Literature Review: The Page--Geilker Experiment,\\
Semiclassical Gravity Inconsistencies,\\
and Constraint-Field Approaches to Quantum Gravity\\[6pt]
\large Cross-Referenced with the DFD Framework\\[6pt]
\normalsize DFD Research Programme}
\author{Cross-Reference Analysis}
\date{23 March 2026}

\begin{document}
\maketitle
\tableofcontents
\newpage

%=============================================================================
\section{Executive Summary}
\label{sec:executive}
%=============================================================================

This review surveys the major theoretical arguments and experimental proposals concerning the quantum vs.\ classical nature of gravity, with particular attention to how DFD's \textbf{quasi-static constraint-field} picture of the scalar field $\psi$ relates to each line of argument.

\textbf{Key finding:} DFD occupies a unique position in the landscape. Its $\psi$ field is \emph{not} a standard semiclassical gravity theory (it does not use $G_{\mu\nu} = 8\pi G \langle \hat{T}_{\mu\nu}\rangle$), nor is it a fully quantized graviton theory. Instead, $\psi$ is a \textbf{constraint field}---analogous to the Coulomb potential $A_0$ in QED---that is determined instantaneously by the matter distribution and carries no independent propagating degrees of freedom. This means:

\begin{enumerate}
\item The Page--Geilker experiment does \emph{not} constrain DFD, because DFD does not use the mean-field semiclassical Einstein equation.
\item The Eppley--Hannah argument does \emph{not} apply, because $\psi$ has no free-particle excitations that could be prepared as low-momentum, short-wavelength probes.
\item The BMV entanglement-witness proposal is \emph{compatible} with DFD: the constraint field can mediate entanglement through its dependence on the full quantum state of matter, not merely its expectation value.
\item The Anastopoulos--Hu critique of Newtonian entanglement arguments \emph{supports} DFD's position: Newtonian gravity (a constraint, not a dynamical degree of freedom) is agnostic to the quantum/classical status of gravitational waves.
\end{enumerate}

%=============================================================================
\part{The Core Papers}
%=============================================================================

%=============================================================================
\section{Page \& Geilker (1981): Indirect Evidence for Quantum Gravity}
\label{sec:page-geilker}
%=============================================================================

\subsection{What They Showed}

\textbf{Reference:} D.N.\ Page and C.D.\ Geilker, ``Indirect Evidence for Quantum Gravity,'' \textit{Phys.\ Rev.\ Lett.}\ \textbf{47}, 979--982 (1981).

Page and Geilker performed a modified Cavendish experiment designed to test the \textbf{semiclassical Einstein equation} (SCE):
\begin{equation}\label{eq:SCE}
G_{\mu\nu} = 8\pi G \langle \hat{T}_{\mu\nu} \rangle,
\end{equation}
in which the classical gravitational field is sourced by the \emph{quantum expectation value} of the stress-energy tensor, rather than by its eigenvalue in a particular branch of the wave function.

\subsubsection{Experimental design}

A radioactive source was used to trigger a quantum branching event: depending on whether the source decayed within a 30-second window, massive lead balls were placed in one of two configurations (A or B) of a Cavendish torsion balance. The quantum event created a superposition $|\Psi\rangle = \alpha |A\rangle + \beta |B\rangle$ with $|\alpha|^2 \approx |\beta|^2 \approx 1/2$.

Under the SCE \eqref{eq:SCE}, the gravitational field should be sourced by $\langle \hat{T}_{\mu\nu}\rangle$, which is the \emph{average} of the two configurations. This average is approximately the same as having no net displacement---the torsion balance should show no deflection.

\subsubsection{Result}

The torsion balance \emph{did} deflect, consistent with the gravitational field being sourced by the \emph{actual} (collapsed/decohered) configuration, not the superposition average. Page and Geilker argued this provided indirect evidence against the SCE and hence indirect evidence \emph{for} quantum gravity.

\subsubsection{Criticisms}

The experiment has been widely criticized:
\begin{itemize}
\item \textbf{Decoherence:} The macroscopic Cavendish apparatus decoheres essentially instantaneously; the relevant question is whether the SCE holds for \emph{microscopic} superpositions, which this experiment cannot probe.
\item \textbf{Interpretation dependence:} The result depends on one's interpretation of quantum mechanics. In Everett (many-worlds), each branch sees a definite configuration; the SCE in each branch gives the correct result. The experiment only constrains the SCE in a \emph{single-world} interpretation where macroscopic superpositions persist.
\item \textbf{Scale:} As Hossenfelder (2012) noted, the experiment operates entirely in the classical regime and does not probe the quantum-gravity interface.
\end{itemize}

\subsection{Relation to DFD}

\textbf{Status: \GREEN{} --- DFD is fully consistent.}

DFD's $\psi$ field is \emph{not} governed by the semiclassical Einstein equation \eqref{eq:SCE}. The DFD field equation in the quasi-static limit is the modified Poisson equation:
\begin{equation}\label{eq:DFD-Poisson}
\nabla \cdot \left[\mu\!\left(\frac{|\nabla\psi|^2}{a_0^2}\right)\nabla\psi\right] = 4\pi G \rho,
\end{equation}
where $\rho$ is the \emph{actual} mass density, not $\langle \hat{\rho} \rangle$. More precisely, in the constraint-field quantization (established i32--i35), $\psi$ is determined by the matter configuration on each branch of the wave function separately. In a Cavendish-type experiment:
\begin{itemize}
\item On branch $|A\rangle$: $\psi$ solves \eqref{eq:DFD-Poisson} with the configuration-A source.
\item On branch $|B\rangle$: $\psi$ solves \eqref{eq:DFD-Poisson} with the configuration-B source.
\item There is \emph{no} branch where $\psi$ is sourced by the average.
\end{itemize}

This is precisely the Coulomb-gauge analogy: in QED, $A_0$ on each branch is $A_0 = \int d^3x'\, \rho(\mathbf{x}')/|\mathbf{x}-\mathbf{x}'|$, determined by the charge configuration on that branch. DFD's $\psi$ behaves identically. The Page--Geilker experiment would show deflection in DFD, exactly as observed.


%=============================================================================
\section{Eppley \& Hannah (1977): The Necessity of Quantizing Gravity}
\label{sec:eppley-hannah}
%=============================================================================

\subsection{What They Argued}

\textbf{Reference:} K.\ Eppley and E.\ Hannah, ``The necessity of quantizing the gravitational field,'' \textit{Found.\ Phys.}\ \textbf{7}, 51--68 (1977).

Eppley and Hannah proposed a thought experiment intended as a \emph{no-go theorem} against classical gravity coupled to quantum matter. Their argument proceeds as follows:

\begin{enumerate}
\item In a classical theory of gravity, gravitational waves need not satisfy $p = \hbar k$ (the de~Broglie relation). One could, in principle, prepare a gravitational wave with short wavelength (high spatial resolution) but arbitrarily small momentum.
\item Such a wave could be used to measure the position of a quantum particle without disturbing its momentum, violating the Heisenberg uncertainty principle.
\item Alternatively, if the measurement \emph{does} disturb the particle, conservation of momentum is violated (since the classical wave carries arbitrarily small momentum).
\item A further variant shows that such interactions would permit superluminal signalling.
\end{enumerate}

The conclusion: gravity must be quantized to avoid these inconsistencies.

\subsection{Subsequent Criticisms}

The Eppley--Hannah argument has been substantially undermined:

\begin{itemize}
\item \textbf{Mattingly (2006):} Showed that the thought experiment requires preparation of gravitational wave states that are physically impossible even in classical GR (the energy density required exceeds black hole formation thresholds). The argument assumes operationalist control over gravity that is not achievable.

\item \textbf{Kent (2018):} Provided a simple refutation: if the gravitational field interacts with matter via the \emph{local quantum state} (not the expectation value), the argument fails entirely. A hybrid theory where gravity responds to the quantum state branch-by-branch (rather than to $\langle\hat{T}_{\mu\nu}\rangle$) evades all three horns of the Eppley--Hannah dilemma. Kent's key point: the alternative considered by Eppley--Hannah (gravity couples to expectation values) is not the only classical option.

\item \textbf{Huggett \& Callender (2001):} Argued that the thought experiment falls short of a no-go theorem because it implicitly assumes that the gravitational wave can be treated as an external classical probe, which begs the question.
\end{itemize}

\subsection{Relation to DFD}

\textbf{Status: \GREEN{} --- DFD evades the argument entirely.}

DFD's $\psi$ has \textbf{no free-particle excitations at the linearized level}. The kinetic Lagrangian $\mathcal{L}_{\rm kin} = \chi - \ln(1+\chi)$ starts at $O(\chi^2)$, i.e., $O((\partial\psi)^4)$. There is no $(\partial\psi)^2$ term, hence no free propagator, hence no $\psi$-quanta that could be prepared as short-wavelength, low-momentum probes. The Eppley--Hannah thought experiment requires gravitational waves with independently tunable wavelength and momentum---this is impossible in DFD because $\psi$ does not propagate freely.

Furthermore, DFD falls precisely into Kent's (2018) escape route: $\psi$ responds to the local matter configuration on each quantum branch, not to $\langle\hat{T}_{\mu\nu}\rangle$. This is the constraint-field quantization: $\psi$ is determined by the matter state via the constraint equation, branch by branch.

The Eppley--Hannah argument constrains theories where gravity has propagating classical degrees of freedom that fail to satisfy quantum commutation relations. DFD's $\psi$ has \emph{no} propagating degrees of freedom (in the free-field limit), so it is outside the scope of the argument.


%=============================================================================
\section{Mattingly (2006): Is Quantum Gravity Necessary?}
\label{sec:mattingly}
%=============================================================================

\subsection{What He Argued}

\textbf{Reference:} J.\ Mattingly, ``Is Quantum Gravity Necessary?'' in \textit{The Universe of General Relativity}, eds.\ A.J.\ Kox and J.\ Eisenstaedt, Einstein Studies Vol.\ 11, pp.\ 327--338 (Birkh\"auser, 2005/2006). See also: J.\ Mattingly, ``Why Eppley and Hannah's thought experiment fails,'' \textit{Phys.\ Rev.\ D} \textbf{73}, 064025 (2006).

Mattingly reviewed the standard arguments for the necessity of quantizing gravity and found them all wanting:

\begin{enumerate}
\item \textbf{The Eppley--Hannah argument} fails because it requires operationalist control over gravitational waves that exceeds what classical GR permits (see Section~\ref{sec:eppley-hannah}).

\item \textbf{The Page--Geilker experiment} tests only the mean-field SCE at macroscopic scales where decoherence has already collapsed the superposition; it does not probe the quantum-gravity interface.

\item \textbf{Formal consistency arguments} (e.g., that hybrid classical-quantum theories are necessarily inconsistent) assume a specific coupling structure that is not the only possibility.

\item \textbf{The measurement problem:} Mattingly argued that the generic solution to the measurement problem tied to semiclassical gravity, combined with the difficulty of quantum gravity alternatives, makes semiclassical gravity a reasonable default option in the absence of decisive experimental evidence.
\end{enumerate}

Mattingly also reviewed a class of recently discovered semiclassical models (stochastic and collapse-based) that bypass the most serious objections, concluding that the question of whether gravity must be quantized is ultimately \emph{empirical}, not settled by theoretical arguments alone.

\subsection{Relation to DFD}

\textbf{Status: \GREEN{} --- DFD is strongly aligned with Mattingly's conclusions.}

Mattingly's central point---that no existing theoretical argument conclusively proves gravity must be quantized---is fully compatible with DFD's framework. DFD goes further than Mattingly's general defense of semiclassical approaches by providing a \emph{specific} mechanism (the constraint field) that:

\begin{itemize}
\item Avoids the mean-field pathology of the standard SCE (because $\psi$ tracks each quantum branch, not $\langle\hat{T}_{\mu\nu}\rangle$).
\item Has no propagating $\psi$-quanta to create Eppley--Hannah-type paradoxes.
\item Is not ``semiclassical gravity'' in the standard sense at all---it is a constraint-field theory where the Newtonian-like potential is a non-dynamical variable determined by matter.
\end{itemize}

DFD provides the concrete realization of the ``third way'' that Mattingly argued should remain open.


%=============================================================================
\section{Carlip (2008): Is Quantum Gravity Necessary?}
\label{sec:carlip}
%=============================================================================

\subsection{What He Argued}

\textbf{Reference:} S.\ Carlip, ``Is Quantum Gravity Necessary?'' \textit{Class.\ Quantum Grav.}\ \textbf{25}, 154010 (2008). arXiv:0803.3456.

Carlip approached the question from a different angle than Mattingly. While acknowledging that theoretical arguments against mixed classical-quantum models are ``strong but not conclusive,'' Carlip emphasized:

\begin{enumerate}
\item The SCE \eqref{eq:SCE} spoils the linearity of quantum mechanics. If gravity is sourced by $\langle\hat{T}_{\mu\nu}\rangle$, the gravitational field becomes a nonlinear functional of the quantum state, which can be used for superluminal signalling (the ``Schr\"odinger cat'' problem of semiclassical gravity).

\item The Page--Geilker experiment, despite its limitations, does show that the mean-field SCE is incompatible with observation at macroscopic scales.

\item The question is ultimately experimental. Carlip reviewed proposals for testing the nonlinearity of the classical-quantum coupling, which could settle whether gravity must be quantized.

\item Carlip was cautiously open to the possibility that a consistent hybrid theory might exist, but emphasized that all known variants (SCE, stochastic semiclassical gravity) have serious problems.
\end{enumerate}

\subsection{Relation to DFD}

\textbf{Status: \GREEN{} --- DFD resolves Carlip's objections.}

Carlip's core concern is that the SCE introduces nonlinearity into quantum mechanics via the $\langle\hat{T}_{\mu\nu}\rangle$ source. DFD \textbf{does not have this problem}:

\begin{itemize}
\item DFD's $\psi$ is a \emph{constraint field} determined by the matter configuration on each branch of the wave function. It does not introduce nonlinearity into the Schr\"odinger equation.
\item The quantum state evolves linearly; $\psi$ is an auxiliary variable that tracks the matter distribution, not a dynamical degree of freedom that feeds back nonlinearly.
\item The superluminal signalling problem of the SCE arises because $\langle\hat{T}_{\mu\nu}\rangle$ is a nonlocal functional of the quantum state. DFD's constraint equation \eqref{eq:DFD-Poisson} is \emph{local}: $\psi$ at point $\mathbf{x}$ depends on the matter distribution through an elliptic equation, just as the Coulomb potential in QED depends on the charge distribution through Poisson's equation. The apparent instantaneity is a gauge artifact (as in Coulomb gauge), not a source of signalling.
\end{itemize}

DFD is precisely the type of theory Carlip suggested \emph{might} exist but had not been constructed: a hybrid approach that avoids both the pathologies of the SCE and the unsolved problems of full quantum gravity.


%=============================================================================
\section{Kafri, Taylor \& Milburn (2014): Classical Channel Model}
\label{sec:kafri}
%=============================================================================

\subsection{What They Argued}

\textbf{Reference:} D.\ Kafri, J.M.\ Taylor, and G.J.\ Milburn, ``A classical channel model for gravitational decoherence,'' \textit{New J.\ Phys.}\ \textbf{16}, 065020 (2014). arXiv:1401.0946.

Kafri, Taylor, and Milburn (KTM) modelled the gravitational interaction between two massive quantum systems as a \emph{classical measurement-and-feedback channel}:

\begin{enumerate}
\item Gravity continuously ``measures'' the position of each mass (extracting classical information).
\item This classical information is used to apply a gravitational force on the other mass.
\item The process introduces decoherence (measurement back-action) equivalent to the Di\'osi--Penrose gravitational decoherence model.
\end{enumerate}

A key result: a classical channel mediating gravity through continuous weak measurement \textbf{cannot generate entanglement} between the two masses (for Gaussian states). This was later used as the basis for the BMV experiment proposal (Section~\ref{sec:bmv}): if entanglement \emph{is} observed, the mediator cannot be a classical channel.

\subsection{Relation to DFD}

\textbf{Status: \YELLOW{} --- Requires careful distinction.}

DFD's $\psi$ is \emph{not} a classical measurement channel in the KTM sense. The crucial difference:

\begin{itemize}
\item \textbf{KTM:} Gravity extracts classical information from the quantum state (a measurement), then feeds it back classically. This is an explicitly information-destroying process that causes decoherence.

\item \textbf{DFD:} The constraint field $\psi$ is determined by the full quantum state of matter on each branch, without any measurement or information extraction. There is no measurement back-action. The constraint equation $\nabla\cdot[\mu\nabla\psi] = 4\pi G\rho$ operates on the matter distribution in each branch coherently.
\end{itemize}

The DFD constraint field is more analogous to the Coulomb interaction in QED, which mediates entanglement perfectly well (electromagnetic interactions routinely entangle charged particles) despite the Coulomb potential being a non-dynamical constraint. The relevant mechanism is not classical information transfer but \emph{constraint-mediated quantum correlation}.

However, the precise mechanism by which DFD's constraint field mediates entanglement requires further theoretical development (see Section~\ref{sec:synthesis}). The key question is whether the $\psi$ constraint equation, applied branch-by-branch, generates the correct entanglement structure. Based on the Coulomb-gauge analogy, the answer should be yes, but a rigorous proof is needed.


%=============================================================================
\section{Tilloy \& Di\'osi (2016): Sourcing Semiclassical Gravity from Spontaneous Localization}
\label{sec:tilloy-diosi}
%=============================================================================

\subsection{What They Argued}

\textbf{Reference:} A.\ Tilloy and L.\ Di\'osi, ``Sourcing semiclassical gravity from spontaneously localized quantum matter,'' \textit{Phys.\ Rev.\ D} \textbf{93}, 024026 (2016). arXiv:1509.08705.

Tilloy and Di\'osi addressed the fundamental inconsistencies of the standard SCE:
\begin{itemize}
\item Schr\"odinger cat sources (a superposition of macroscopically distinct mass configurations sources a gravitational field that is the average---physically absurd).
\item Superluminal signalling (the nonlinearity of $\langle\hat{T}_{\mu\nu}\rangle$ as a source allows signal propagation outside the light cone).
\item Violation of the Born rule.
\end{itemize}

Their resolution: replace the standard quantum state with a \emph{spontaneously localized} quantum state (using GRW or CSL collapse models). The gravitational field is then sourced by the \emph{collapsed} matter distribution, which is always approximately localized. This gives:
\begin{itemize}
\item A stochastic semiclassical gravity theory in the Newtonian limit.
\item Recovery of the Newtonian pair potential (no single-particle self-interaction).
\item An additional gravitational decoherence term that depends on the specific collapse model.
\item Quantitative, testable predictions that distinguish it from fully quantized gravity.
\end{itemize}

\subsection{Relation to DFD}

\textbf{Status: \GREEN{} --- DFD achieves the same goals through a different mechanism.}

DFD and Tilloy--Di\'osi share the same diagnosis of the problem (the SCE is inconsistent due to mean-field sourcing) and arrive at structurally similar solutions (the gravitational field tracks the actual matter configuration, not the expectation value). However, the mechanisms differ:

\begin{itemize}
\item \textbf{Tilloy--Di\'osi:} Invokes an explicit collapse mechanism (GRW/CSL) to ensure the matter state is always approximately localized. The gravitational field then tracks the collapsed state.

\item \textbf{DFD:} Does not require a specific collapse mechanism. The constraint-field structure of $\psi$ naturally ensures that on each branch of the wave function, $\psi$ is determined by the matter configuration on \emph{that branch}. No collapse is needed; the constraint equation operates on whatever matter state is present.
\end{itemize}

DFD is therefore more general than Tilloy--Di\'osi: it does not commit to a specific resolution of the measurement problem. Whether the wave function collapses (GRW), branches (Everett), or is guided (de~Broglie--Bohm), the DFD constraint equation produces the correct $\psi$ for the relevant matter configuration.

An important shared prediction: \textbf{no single-particle gravitational self-interaction}. In both DFD and Tilloy--Di\'osi, the gravitational potential of a particle does not act back on that same particle (this would be the Schr\"odinger--Newton self-energy, which Anastopoulos \& Hu (2014) showed is problematic). DFD achieves this naturally because $\psi$ mediates interactions between \emph{different} matter sources.


%=============================================================================
\section{Anastopoulos \& Hu (2014--2022): Problems with Semiclassical Gravity}
\label{sec:anastopoulos-hu}
%=============================================================================

\subsection{What They Argued}

\textbf{Key References:}
\begin{itemize}
\item C.\ Anastopoulos and B.L.\ Hu, ``Problems with the Newton--Schr\"odinger Equations,'' \textit{New J.\ Phys.}\ \textbf{16}, 085007 (2014). arXiv:1403.4921.
\item C.\ Anastopoulos and B.L.\ Hu, ``Gravity, Quantum Fields and Quantum Information: Problems with Classical Channel and Stochastic Theories,'' \textit{Entropy} \textbf{24}, 490 (2022).
\item C.\ Anastopoulos and B.L.\ Hu, ``Comment on `A Spin Entanglement Witness for Quantum Gravity'\thinspace'' (2018). arXiv:1804.11315.
\end{itemize}

Anastopoulos and Hu have been among the most careful analysts of the semiclassical gravity landscape. Their key results:

\subsubsection{Problems with the Newton--Schr\"odinger equations (2014)}

The Newton--Schr\"odinger equation (NSE):
\begin{equation}\label{eq:NSE}
i\hbar\frac{\partial\Psi}{\partial t} = \left[-\frac{\hbar^2}{2m}\nabla^2 + V_{\rm ext} - Gm^2\int d^3x'\,\frac{|\Psi(\mathbf{x}',t)|^2}{|\mathbf{x}-\mathbf{x}'|}\right]\Psi
\end{equation}
is often claimed to be the non-relativistic limit of the SCE. Anastopoulos and Hu showed this is \textbf{false}:
\begin{itemize}
\item The NSE is a single-particle equation; the SCE is a field equation. They do not have the same limit.
\item The wave function in the NSE makes sense only as a mean-field description of a system of $N$ particles as $N \to \infty$, not as the wave function of a single particle.
\item The NSE predicts gravitational self-interaction of a single particle (the last term in \eqref{eq:NSE}), which has no physical basis from GR + QFT.
\end{itemize}

\subsubsection{Critique of BMV-type arguments (2018, 2022)}

Anastopoulos and Hu argued that gravity-induced entanglement through Newtonian forces is \textbf{agnostic to the quantum or classical nature of the gravitational true degrees of freedom} (i.e., gravitational waves). Their reasoning:

\begin{itemize}
\item In GR, the Newtonian potential is not a dynamical degree of freedom---it is a \emph{constraint} determined by the Hamiltonian constraint equation.
\item The true dynamical degrees of freedom of gravity are the transverse-traceless gravitational waves.
\item Entanglement mediated by the Newtonian potential (a constraint) does not prove that the dynamical degrees of freedom (gravitational waves) are quantum.
\item The Anastopoulos--Blencowe--Hu model of gravitational decoherence, which involves classical/stochastic gravitational waves plus quantum Newtonian constraints, is fully compatible with gravity-induced entanglement in the Newtonian regime.
\end{itemize}

\subsection{Relation to DFD}

\textbf{Status: \GREEN{} --- DFD is strongly aligned with Anastopoulos--Hu.}

The Anastopoulos--Hu analysis provides \emph{direct theoretical support} for DFD's constraint-field picture:

\begin{enumerate}
\item \textbf{Constraint vs.\ dynamical degrees of freedom:} Anastopoulos and Hu's distinction between the Newtonian potential (constraint) and gravitational waves (dynamical) maps directly onto DFD's framework. DFD's $\psi$ is the analogue of the Newtonian constraint. DFD does not quantize $\psi$ as an independent degree of freedom, consistent with its constraint status.

\item \textbf{NSE critique:} DFD's field equation \eqref{eq:DFD-Poisson} does \emph{not} reduce to the NSE. It is a constraint equation sourced by $\rho$ (not $|\Psi|^2$), and it operates at the multi-particle level as a field equation, not as a single-particle self-interaction. DFD naturally avoids the NSE pathologies that Anastopoulos and Hu identified.

\item \textbf{Entanglement mediation:} The Anastopoulos--Hu observation that Newtonian-constraint-mediated entanglement does not imply quantum gravitational waves is precisely the DFD position. In DFD, $\psi$-mediated entanglement between matter systems would be expected (just as Coulomb-mediated entanglement occurs in QED), without implying that $\psi$ is itself a quantum field with its own Fock space.

\item \textbf{Gravitational decoherence:} DFD predicts no fundamental gravitational decoherence from the $\psi$ sector (since $\psi$ is a constraint, not a stochastic classical field). Any decoherence would come from the gravitational-wave sector, which DFD inherits from GR. This is consistent with the Anastopoulos--Blencowe--Hu model.
\end{enumerate}


%=============================================================================
\section{Bose, Marletto, Vedral (2017): The BMV Experiment}
\label{sec:bmv}
%=============================================================================

\subsection{What They Proposed}

\textbf{References:}
\begin{itemize}
\item S.\ Bose \textit{et al.}, ``Spin Entanglement Witness for Quantum Gravity,'' \textit{Phys.\ Rev.\ Lett.}\ \textbf{119}, 240401 (2017). arXiv:1707.06036.
\item C.\ Marletto and V.\ Vedral, ``Gravitationally Induced Entanglement between Two Massive Particles is Sufficient Evidence of Quantum Effects in Gravity,'' \textit{Phys.\ Rev.\ Lett.}\ \textbf{119}, 240402 (2017).
\end{itemize}

The BMV (Bose--Marletto--Vedral) experiment proposes to test whether gravity can mediate entanglement between two mesoscopic masses:

\begin{enumerate}
\item Two micron-scale masses are each placed in a spatial superposition using Stern--Gerlach-type splitting (spin-dependent forces).
\item The masses interact gravitationally while in superposition.
\item After recombination, spin-spin correlations are measured.
\item If the masses become entangled, the gravitational mediator must be ``non-classical'' (in the sense of not being a local operations and classical communication, or LOCC, channel).
\end{enumerate}

The theoretical basis is the LOCC theorem: two quantum systems can become entangled through the mediation of a third system only if that mediator is capable of transmitting quantum information. A purely classical channel (like KTM's model) cannot generate entanglement.

\subsection{Relation to DFD}

\textbf{Status: \YELLOW{} --- DFD predicts entanglement, but the interpretation differs.}

DFD predicts that the BMV experiment \emph{will} show entanglement, but this does \emph{not} imply that gravity is quantized in the conventional sense:

\begin{itemize}
\item \textbf{DFD's constraint field mediates entanglement:} Just as the Coulomb potential in QED mediates entanglement between charged particles (without anyone claiming that $A_0$ is a quantum field), DFD's $\psi$ constraint field mediates entanglement between massive particles. The mechanism is the branch-dependent constraint: on each branch, $\psi$ is determined by the matter configuration, creating correlations between the two masses.

\item \textbf{The LOCC theorem does not apply as stated:} The BMV argument assumes gravity is either (a)~a quantum mediator or (b)~a classical LOCC channel. DFD's constraint field is \emph{neither}---it is a constraint determined by the quantum state of matter, not an independent mediator with its own quantum or classical degrees of freedom. The LOCC theorem requires an independent mediating system; a constraint is not an independent system.

\item \textbf{Anastopoulos--Hu's critique applies:} As discussed in Section~\ref{sec:anastopoulos-hu}, entanglement through the Newtonian constraint does not prove the dynamical degrees of freedom of gravity are quantum. DFD would predict BMV entanglement from the constraint sector while remaining agnostic about (or even classical regarding) gravitational-wave quantization.

\item \textbf{Recent development---Aziz \& Howl (2025):} The Nature paper ``Classical theories of gravity produce entanglement'' showed that when matter is described by full QFT (rather than non-relativistic quantum mechanics), local classical gravitational theories can generate entanglement through physical, local processes. This provides additional theoretical support for DFD's position. However, this result is contested (see, e.g., the Newton--Cartan analysis by Christodoulou et al., 2025).
\end{itemize}

\textbf{Experimental prediction:} DFD predicts the BMV experiment will observe entanglement, with the entanglement phase given by the standard Newtonian result $\phi = Gm^2 t/(\hbar d)$ (where $d$ is the separation). In the DFD MOND regime ($a \lesssim a_0$), there would be a \emph{modified} entanglement phase due to the $\mu$-function enhancement, but this regime is far below the sensitivity of current BMV proposals.


%=============================================================================
\section{Belenchia, Wald, et al.\ (2018): Quantum Superposition and Gravitational Entanglement}
\label{sec:belenchia-wald}
%=============================================================================

\subsection{What They Argued}

\textbf{Reference:} A.\ Belenchia, R.M.\ Wald, F.\ Giacomini, E.\ Castro-Ruiz, \v{C}.\ Brukner, and M.\ Aspelmeyer, ``Quantum superposition of massive objects and the quantization of gravity,'' \textit{Phys.\ Rev.\ D} \textbf{98}, 126009 (2018). arXiv:1807.07015.

Belenchia et al.\ analyzed a Gedankenexperiment involving quantum superpositions of charged and/or massive bodies. Their key insight:

\begin{itemize}
\item A quantum body in a spatial superposition must be viewed as \textbf{entangled with its own Newtonian-like gravitational field}. The ``field'' in each branch takes the value determined by the position of the mass in that branch.
\item This body--field entanglement is essential for understanding how the body may become entangled with \emph{other} massive bodies via gravitational interactions.
\item They corrected earlier analyses that failed to account for conservation of the center of mass of the isolated system.
\end{itemize}

\subsection{Relation to DFD}

\textbf{Status: \GREEN{} --- DFD naturally implements this picture.}

The Belenchia--Wald picture of body--field entanglement is \emph{exactly} what DFD's constraint-field quantization produces:

\begin{itemize}
\item On branch $|\mathbf{x}_1\rangle$ (mass at position $\mathbf{x}_1$): $\psi$ solves the constraint equation with source at $\mathbf{x}_1$.
\item On branch $|\mathbf{x}_2\rangle$ (mass at position $\mathbf{x}_2$): $\psi$ solves the constraint equation with source at $\mathbf{x}_2$.
\item The total state is: $|\Psi\rangle = \alpha|\mathbf{x}_1\rangle|\psi_1\rangle + \beta|\mathbf{x}_2\rangle|\psi_2\rangle$, where $|\psi_i\rangle$ denotes the field configuration determined by the constraint.
\end{itemize}

This is body--field entanglement through the constraint, without $\psi$ being an independent quantum degree of freedom. The constraint \emph{creates} the entanglement by correlating the field configuration with the matter state on each branch.


%=============================================================================
\section{Kent (2018): Simple Refutation of Eppley--Hannah}
\label{sec:kent}
%=============================================================================

\subsection{What He Argued}

\textbf{Reference:} A.\ Kent, ``Simple Refutation of the Eppley--Hannah Argument,'' \textit{Class.\ Quantum Grav.}\ \textbf{35}, 245008 (2018). arXiv:1807.08708.

Kent's refutation is concise and powerful: if the gravitational field interacts with matter via the \emph{local quantum state} (rather than the expectation value $\langle\hat{T}_{\mu\nu}\rangle$), the Eppley--Hannah argument fails. Specifically:

\begin{itemize}
\item Eppley--Hannah considered only two options: (a)~gravity is quantized, or (b)~gravity couples to $\langle\hat{T}_{\mu\nu}\rangle$.
\item Kent identified a third option: gravity is ``classical'' in the sense of not having its own quantum fluctuations, but it couples to the local quantum state of matter (not the expectation value). In this case, there are no independently tunable classical gravitational waves, and the measurement-disturbance argument fails.
\item Kent's option (c) is plausibly natural for a hybrid theory where the gravitational field is determined by the matter configuration on each branch.
\end{itemize}

\subsection{Relation to DFD}

\textbf{Status: \GREEN{} --- DFD is the concrete realization of Kent's option (c).}

Kent's ``option (c)'' is precisely what DFD implements:
\begin{itemize}
\item $\psi$ has no independent quantum fluctuations (no $\psi$-quanta).
\item $\psi$ is determined by the matter distribution on each branch via the constraint equation.
\item There are no classical $\psi$-waves that could be independently prepared and used as probes.
\end{itemize}

DFD provides the explicit Lagrangian, field equation, and quantization scheme for the type of theory Kent argued should be viable.


%=============================================================================
\section{Oppenheim (2023): A Postquantum Theory of Classical Gravity}
\label{sec:oppenheim}
%=============================================================================

\subsection{What He Argued}

\textbf{Reference:} J.\ Oppenheim, ``A Postquantum Theory of Classical Gravity?'' \textit{Phys.\ Rev.\ X} \textbf{13}, 041040 (2023). arXiv:1811.03116.

Oppenheim constructed the most rigorous known theory of \emph{fundamentally classical gravity coupled to quantum matter}:

\begin{enumerate}
\item The theory is hybrid classical-quantum: the metric lives on a classical phase space, while matter lives in a Hilbert space.
\item The coupling is \textbf{fundamentally stochastic}: both the metric and quantum fields experience stochastic fluctuations.
\item The dynamics is linear in the density matrix, completely positive, and trace preserving (CPTP)---avoiding the pathologies of the SCE.
\item The theory reduces to Einstein's GR in the classical limit.
\item The stochastic coupling preserves the uncertainty principle for quantum matter and prevents superluminal signalling.
\item The price: fundamental decoherence and diffusion in the gravitational degrees of freedom.
\end{enumerate}

Key predictions distinguishing Oppenheim's theory from quantum gravity:
\begin{itemize}
\item Anomalous heating of matter due to gravitational diffusion.
\item Decoherence of spatial superpositions at a rate determined by the gravitational coupling.
\item A trade-off between decoherence and diffusion that is experimentally testable.
\end{itemize}

\subsection{Relation to DFD}

\textbf{Status: \YELLOW{} --- DFD and Oppenheim share motivations but differ structurally.}

Similarities:
\begin{itemize}
\item Both avoid the mean-field SCE and its pathologies.
\item Both treat the gravitational sector as not fully quantized.
\item Both achieve CPTP dynamics for the matter sector.
\end{itemize}

Differences:
\begin{itemize}
\item \textbf{Stochasticity:} Oppenheim's theory is fundamentally stochastic; the classical-quantum coupling introduces irreducible noise. DFD's constraint-field coupling is \textbf{deterministic} on each branch---$\psi$ is uniquely determined by the matter configuration. DFD has no fundamental gravitational noise or diffusion from the $\psi$ sector.

\item \textbf{Decoherence:} Oppenheim predicts fundamental gravitational decoherence proportional to the mass separation. DFD predicts \textbf{no} $\psi$-induced decoherence (since $\psi$ is a constraint, not a stochastic variable). This is a sharp experimental discriminator: if Oppenheim-type decoherence is observed, it disfavors DFD's deterministic constraint picture (though DFD's gravitational-wave sector could still contribute decoherence).

\item \textbf{Independent degrees of freedom:} Oppenheim's metric has independent classical degrees of freedom (gravitational waves are classical). DFD's $\psi$ has \emph{no} independent degrees of freedom at the linearized level. The DFD gravitational-wave sector is inherited from GR and may be quantized or classical.

\item \textbf{Testability:} Oppenheim's theory makes specific predictions for anomalous heating rates. DFD's predictions in this sector are different: no anomalous heating from $\psi$, but MOND-scale modifications to entanglement phases.
\end{itemize}


%=============================================================================
\part{Constraint-Field and Coulomb-Gauge Approaches}
%=============================================================================

%=============================================================================
\section{The Coulomb-Gauge Analogy for Gravity}
\label{sec:coulomb-gauge}
%=============================================================================

\subsection{The Analogy in Electromagnetism}

In the Coulomb gauge ($\nabla \cdot \mathbf{A} = 0$), the electromagnetic potential decomposes into:
\begin{itemize}
\item $A_0$ (the Coulomb potential): determined \emph{instantaneously} by the charge distribution via $\nabla^2 A_0 = -\rho/\epsilon_0$. This is a \textbf{constraint}, not a dynamical variable. It has no conjugate momentum ($\pi_0 = \partial\mathcal{L}/\partial\dot{A}_0 = 0$) and no associated quanta.
\item $\mathbf{A}_T$ (transverse vector potential): the dynamical degree of freedom, satisfying a wave equation. Quantization of $\mathbf{A}_T$ produces photons.
\end{itemize}

Despite being a ``classical'' constraint, $A_0$ mediates the Coulomb interaction that entangles charged particles:
\begin{equation}
V_{\rm Coulomb} = \frac{1}{2}\int d^3x\,d^3x'\,\frac{\rho(\mathbf{x})\rho(\mathbf{x}')}{4\pi\epsilon_0|\mathbf{x}-\mathbf{x}'|}.
\end{equation}
This interaction generates entanglement in every atom, every molecule, every solid---without $A_0$ ever being ``quantized'' as an independent field. The entanglement arises because $A_0$ creates a correlation between the quantum states of different charged particles through the constraint.

The apparent instantaneity of $A_0$ does not violate causality. Observable quantities (electromagnetic fields, particle trajectories) propagate causally. The ``instantaneous'' propagation of $A_0$ is a gauge artifact---it is cancelled by corresponding terms in $\mathbf{A}_T$ when computing physical observables.

\subsection{The Gravitational Analogue: Poisson Gauge}

In the ADM decomposition of general relativity, the analogous decomposition is:
\begin{itemize}
\item \textbf{Lapse $N$ and shift $N^i$:} Non-dynamical (Lagrange multipliers for the Hamiltonian and momentum constraints). Analogous to $A_0$.
\item \textbf{Transverse-traceless spatial metric perturbations $h_{ij}^{TT}$:} The dynamical degrees of freedom (gravitational waves). Analogous to $\mathbf{A}_T$.
\end{itemize}

The Newtonian potential $\Phi$ emerges from the Hamiltonian constraint (analogous to Gauss's law):
\begin{equation}
\nabla^2 \Phi = 4\pi G \rho.
\end{equation}
This is an \emph{elliptic} equation---no time derivatives---determining $\Phi$ instantaneously from $\rho$. In the Poisson gauge (the gravitational analogue of Coulomb gauge), the scalar and vector gravitational potentials depend solely on the instantaneous stress-energy distribution.

As noted by Bertschinger (2002), this apparent action-at-a-distance is a gauge artifact, just as in Coulomb-gauge electromagnetism.

\subsection{DFD's $\psi$ as a Generalized Newtonian Constraint}

DFD's scalar field $\psi$ is a \emph{nonlinear generalization} of the Newtonian potential:
\begin{equation}
\nabla\cdot[\mu(\chi)\nabla\psi] = 4\pi G\rho, \qquad \chi = |\nabla\psi|^2/a_0^2.
\end{equation}

This is still an elliptic constraint equation (no time derivatives of $\psi$ appear in the quasi-static limit). The $\mu$-function introduces MOND phenomenology without changing the constraint character of the equation.

The key structural features inherited from the Coulomb/Poisson-gauge analogy:
\begin{enumerate}
\item $\psi$ has no conjugate momentum: $\pi_\psi = 0$ (primary constraint in the Dirac--Bergmann sense).
\item $\psi$ is determined at each instant by the matter distribution.
\item $\psi$ mediates interactions (including entanglement) between matter systems without being an independent quantum degree of freedom.
\item The apparent instantaneity is a gauge artifact; physical observables propagate causally.
\end{enumerate}


%=============================================================================
\section{Constraint Quantization in Scalar-Tensor Gravity}
\label{sec:constraint-quantization}
%=============================================================================

\subsection{ADM Formalism for Scalar-Tensor Theories}

In the ADM (Arnowitt--Deser--Misner) formalism, spacetime is foliated into spatial hypersurfaces $\Sigma_t$, and the dynamics is formulated in terms of the spatial metric $h_{ij}$, its conjugate momentum $\pi^{ij}$, plus any additional fields and their momenta.

For a scalar-tensor theory with scalar field $\psi$, the canonical analysis yields:
\begin{itemize}
\item \textbf{Hamiltonian constraint:} $\mathcal{H} \approx 0$ (analogous to $\nabla\cdot\mathbf{E} = \rho$ in QED).
\item \textbf{Momentum constraints:} $\mathcal{H}_i \approx 0$.
\item \textbf{Scalar field constraint:} If the scalar field Lagrangian has no $\dot{\psi}^2$ term (as in DFD's $\mathcal{L}_{\rm kin} \sim (\partial\psi)^4$), then $\pi_\psi = 0$ is a \emph{primary constraint}.
\end{itemize}

The Dirac--Bergmann constraint analysis then proceeds:
\begin{enumerate}
\item $\pi_\psi \approx 0$ is a primary constraint.
\item Consistency ($\dot{\pi}_\psi \approx 0$) generates a secondary constraint: the field equation for $\psi$.
\item Together, these constraints eliminate $\psi$ and $\pi_\psi$ from the independent phase-space variables.
\item The Dirac bracket of the remaining variables accounts for the $\psi$-mediated interaction.
\end{enumerate}

This is precisely the procedure used in QED to eliminate $A_0$ and its momentum, leaving only the transverse photon degrees of freedom in the physical Hilbert space.

\subsection{The Wheeler--DeWitt Equation and DFD}

In canonical quantum gravity, the Hamiltonian constraint becomes the Wheeler--DeWitt (WDW) equation:
\begin{equation}
\hat{\mathcal{H}}|\Psi\rangle = 0.
\end{equation}

For DFD, the $\psi$-sector constraint is simpler: it is an elliptic equation (not a hyperbolic wave equation), and its quantum implementation is straightforward:
\begin{equation}
\hat{\nabla}\cdot[\mu(\hat{\chi})\hat{\nabla}\hat{\psi}] = 4\pi G\hat{\rho},
\end{equation}
which determines $\hat{\psi}$ in terms of $\hat{\rho}$ at each instant. This is the operator analogue of Gauss's law $\hat{\nabla}\cdot\hat{\mathbf{E}} = \hat{\rho}$ in QED.

The physical Hilbert space is the space of states satisfying this constraint. In practice, this means $\psi$ is not an independent operator---it is a functional of the matter operators, determined by the constraint.


%=============================================================================
\part{Synthesis and DFD Assessment}
\label{sec:synthesis}
%=============================================================================

%=============================================================================
\section{The Landscape of Approaches}
%=============================================================================

\begin{table}[htbp]
\centering
\small
\begin{tabular}{@{}p{3.2cm}p{2.8cm}p{2.8cm}p{2.8cm}p{2.5cm}@{}}
\toprule
\textbf{Theory} & \textbf{Gravity sector} & \textbf{Avoids SCE pathology?} & \textbf{Predicts BMV entanglement?} & \textbf{Fundamental decoherence?} \\
\midrule
Standard SCE & Mean-field classical & No & Unclear & No \\
KTM classical channel & Classical channel & Yes & No & Yes \\
Tilloy--Di\'osi (GRW) & Stochastic semiclassical & Yes & Yes (reduced) & Yes \\
Oppenheim (2023) & Stochastic classical & Yes & Depends on parameters & Yes \\
Full quantum gravity & Fully quantized & Yes & Yes & No (fundamental) \\
\textbf{DFD constraint field} & \textbf{Constraint (non-dynamical)} & \textbf{Yes} & \textbf{Yes} & \textbf{No (from $\psi$)} \\
\bottomrule
\end{tabular}
\caption{Comparison of approaches to the quantum/classical gravity interface. DFD's constraint-field picture occupies a unique position: it avoids SCE pathologies, predicts BMV entanglement, and introduces no fundamental decoherence from the $\psi$ sector.}
\label{tab:comparison}
\end{table}


%=============================================================================
\section{What DFD Gets Right}
\label{sec:dfd-strengths}
%=============================================================================

\begin{enumerate}
\item \textbf{Avoids the Page--Geilker problem:} $\psi$ tracks each branch, not the expectation value. The Cavendish torsion balance deflects correctly.

\item \textbf{Evades Eppley--Hannah:} No propagating $\psi$-quanta to create the paradox. Implements Kent's ``option (c)'' explicitly.

\item \textbf{Consistent with Mattingly's analysis:} DFD provides the concrete ``third way'' between full quantization and the pathological SCE.

\item \textbf{Resolves Carlip's concerns:} No nonlinear feedback into the Schr\"odinger equation; no superluminal signalling.

\item \textbf{Aligned with Anastopoulos--Hu:} The distinction between constraint (Newtonian potential / $\psi$) and dynamical degrees of freedom (gravitational waves) is built into DFD's structure.

\item \textbf{Naturally implements Belenchia--Wald body-field entanglement:} The constraint equation creates correlated body-field states on each branch.

\item \textbf{Avoids NSE pathologies:} No single-particle gravitational self-interaction; the constraint equation is a multi-body field equation.

\item \textbf{The Coulomb-gauge analogy is exact at the structural level:} $\psi$ is to gravity what $A_0$ is to electromagnetism---a constraint field with no independent quanta that nevertheless mediates interactions and entanglement.
\end{enumerate}


%=============================================================================
\section{Open Questions and Future Work}
\label{sec:open-questions}
%=============================================================================

\begin{enumerate}
\item \textbf{Rigorous entanglement proof:} While the Coulomb-gauge analogy strongly suggests that DFD's constraint field mediates entanglement, a rigorous calculation showing that two masses interacting through $\psi$ become entangled (with the correct entanglement entropy) has not yet been performed within the DFD formalism. This is the top priority.

\item \textbf{MOND-regime entanglement signatures:} DFD predicts modified gravitational interactions in the $a \lesssim a_0$ regime. Does this lead to anomalous entanglement phases in a hypothetical ultra-sensitive BMV experiment? The MOND $\mu$-function enhancement could produce a distinctive signature.

\item \textbf{Gravitational-wave sector:} DFD's constraint-field analysis covers only the $\psi$ (scalar/Newtonian) sector. The gravitational-wave sector is inherited from GR. Whether these transverse-traceless modes should be quantized (producing gravitons) or treated classically/stochastically (as in Anastopoulos--Blencowe--Hu or Oppenheim) is an independent question that DFD does not yet answer.

\item \textbf{Oppenheim-type decoherence tests:} Oppenheim's theory predicts specific decoherence and diffusion rates. DFD predicts \emph{no} $\psi$-induced decoherence. Experimental measurement of gravitational decoherence rates would distinguish the two approaches. If Oppenheim-type decoherence is observed at the predicted rates, it would indicate that the gravitational sector has stochastic classical dynamics beyond what DFD's deterministic constraint provides.

\item \textbf{Aziz \& Howl (2025) and the entanglement debate:} The recent Nature paper claiming that classical gravity can produce entanglement when matter is treated in full QFT supports DFD's position but is contested. Resolving this debate (likely through the distinction between constraint-mediated and channel-mediated entanglement) would strengthen DFD's theoretical foundation.

\item \textbf{Post-Newtonian corrections:} The constraint-field analysis is exact in the quasi-static (Newtonian) limit. At post-Newtonian order, $\psi$ acquires time-derivative terms from the relativistic k-essence Lagrangian. The constraint character may be modified. A full post-Newtonian analysis of DFD's quantization is needed.

\item \textbf{Connection to the spectral action and $\alpha$:} DFD derives the fine-structure constant from the spectral action on $\mathbb{C}P^2 \times S^3$. The constraint-field status of $\psi$ should be derivable from the spectral geometry. Establishing this connection would unify the $\alpha$ derivation with the quantum gravity analysis.
\end{enumerate}


%=============================================================================
\section{Conclusion}
\label{sec:conclusion}
%=============================================================================

The literature on semiclassical gravity and the quantum/classical nature of the gravitational interaction has converged on a clear picture:

\begin{enumerate}
\item \textbf{The standard SCE is untenable} as a fundamental theory (Page--Geilker, Carlip, Anastopoulos--Hu).
\item \textbf{No theoretical argument conclusively proves gravity must be quantized} (Mattingly, Kent).
\item \textbf{A constraint-field approach} (where the Newtonian potential is a non-dynamical variable determined by the matter distribution on each quantum branch) \textbf{evades all known objections} to classical/hybrid gravity theories.
\item \textbf{Entanglement through constraints} is a well-established phenomenon in QED (Coulomb interaction) and does not require the mediating field to be independently quantized.
\end{enumerate}

DFD's scalar field $\psi$ is precisely such a constraint field. Its quasi-static, non-propagating character at the linearized level places it outside the scope of the Eppley--Hannah argument, the Page--Geilker test, and the KTM classical-channel no-go theorem. At the same time, it predicts the entanglement signatures that the BMV experiment seeks to observe.

This positions DFD uniquely in the landscape: it is neither standard semiclassical gravity (which is inconsistent) nor full quantum gravity (which remains unsolved). It is a \textbf{constraint-field theory} that implements the gravitational analogue of the Coulomb gauge in QED, mediating interactions and entanglement between quantum matter systems without introducing independent gravitational quanta for the scalar sector.

The next steps are: (1) rigorous proof that the DFD constraint field generates the correct entanglement structure, (2) analysis of MOND-regime entanglement signatures as a distinctive DFD prediction, and (3) connection to the spectral-action derivation of $\alpha$ to unify the theoretical framework.


%=============================================================================
\section*{Bibliography}
\label{sec:bibliography}
%=============================================================================

\begin{enumerate}
\item D.N.\ Page and C.D.\ Geilker, ``Indirect Evidence for Quantum Gravity,'' \textit{Phys.\ Rev.\ Lett.}\ \textbf{47}, 979--982 (1981).

\item K.\ Eppley and E.\ Hannah, ``The necessity of quantizing the gravitational field,'' \textit{Found.\ Phys.}\ \textbf{7}, 51--68 (1977).

\item J.\ Mattingly, ``Is Quantum Gravity Necessary?'' in \textit{The Universe of General Relativity}, Einstein Studies Vol.\ 11, pp.\ 327--338 (Birkh\"auser, 2006).

\item J.\ Mattingly, ``Why Eppley and Hannah's thought experiment fails,'' \textit{Phys.\ Rev.\ D} \textbf{73}, 064025 (2006).

\item S.\ Carlip, ``Is Quantum Gravity Necessary?'' \textit{Class.\ Quantum Grav.}\ \textbf{25}, 154010 (2008). arXiv:0803.3456.

\item D.\ Kafri, J.M.\ Taylor, and G.J.\ Milburn, ``A classical channel model for gravitational decoherence,'' \textit{New J.\ Phys.}\ \textbf{16}, 065020 (2014). arXiv:1401.0946.

\item A.\ Tilloy and L.\ Di\'osi, ``Sourcing semiclassical gravity from spontaneously localized quantum matter,'' \textit{Phys.\ Rev.\ D} \textbf{93}, 024026 (2016). arXiv:1509.08705.

\item C.\ Anastopoulos and B.L.\ Hu, ``Problems with the Newton--Schr\"odinger Equations,'' \textit{New J.\ Phys.}\ \textbf{16}, 085007 (2014). arXiv:1403.4921.

\item C.\ Anastopoulos and B.L.\ Hu, ``Gravity, Quantum Fields and Quantum Information: Problems with Classical Channel and Stochastic Theories,'' \textit{Entropy} \textbf{24}, 490 (2022).

\item C.\ Anastopoulos and B.L.\ Hu, ``Comment on `A Spin Entanglement Witness for Quantum Gravity'\thinspace'' (2018). arXiv:1804.11315.

\item S.\ Bose \textit{et al.}, ``Spin Entanglement Witness for Quantum Gravity,'' \textit{Phys.\ Rev.\ Lett.}\ \textbf{119}, 240401 (2017). arXiv:1707.06036.

\item C.\ Marletto and V.\ Vedral, ``Gravitationally Induced Entanglement between Two Massive Particles is Sufficient Evidence of Quantum Effects in Gravity,'' \textit{Phys.\ Rev.\ Lett.}\ \textbf{119}, 240402 (2017).

\item A.\ Belenchia, R.M.\ Wald, F.\ Giacomini, E.\ Castro-Ruiz, \v{C}.\ Brukner, and M.\ Aspelmeyer, ``Quantum superposition of massive objects and the quantization of gravity,'' \textit{Phys.\ Rev.\ D} \textbf{98}, 126009 (2018). arXiv:1807.07015.

\item A.\ Kent, ``Simple Refutation of the Eppley--Hannah Argument,'' \textit{Class.\ Quantum Grav.}\ \textbf{35}, 245008 (2018). arXiv:1807.08708.

\item J.\ Oppenheim, ``A Postquantum Theory of Classical Gravity?'' \textit{Phys.\ Rev.\ X} \textbf{13}, 041040 (2023). arXiv:1811.03116.

\item M.\ Bahrami, A.\ Gro\ss ardt, S.\ Donadi, and A.\ Bassi, ``The Schr\"odinger--Newton equation and its foundations,'' \textit{New J.\ Phys.}\ \textbf{16}, 115007 (2014). arXiv:1407.4370.

\item M.\ Derakhshani, ``Newtonian Semiclassical Gravity in the Ghirardi--Rimini--Weber Theory with Matter Density Ontology,'' \textit{Phys.\ Lett.\ A} \textbf{378}, 990--998 (2014). arXiv:1304.0471.

\item A.\ Tilloy and L.\ Di\'osi, ``Principle of Least Decoherence for Newtonian Semiclassical Gravity,'' \textit{Phys.\ Rev.\ D} \textbf{96}, 104045 (2017).

\item N.\ Huggett and C.\ Callender, ``Why Quantize Gravity (or Any Other Field for That Matter)?'' \textit{Philos.\ Sci.}\ \textbf{68}, S382--S394 (2001).

\item S.\ Hossenfelder, ``A real thought experiment that shows virtually nothing'' (Backreaction blog, 2012).

\item A.\ Aziz and R.\ Howl, ``Classical theories of gravity produce entanglement,'' \textit{Nature} (2025). arXiv:2510.19714.

\item E.\ Bertschinger, ``Cosmological Dynamics'' (MIT lecture notes, 2002). Discussion of Poisson gauge.

\item R.\ Arnowitt, S.\ Deser, and C.W.\ Misner, ``The dynamics of general relativity,'' \textit{Gen.\ Rel.\ Grav.}\ \textbf{40}, 1997--2027 (2008). (Reprint of 1962 article.)

\item B.S.\ DeWitt, ``Quantum Theory of Gravity. I. The Canonical Theory,'' \textit{Phys.\ Rev.}\ \textbf{160}, 1113--1148 (1967).

\item P.A.M.\ Dirac, \textit{Lectures on Quantum Mechanics} (Belfer Graduate School, 1964).
\end{enumerate}


\end{document}
